% Improved LaTeX Document for A4 Paper - Lean Manufacturing Report
% Optimized for Overleaf

\documentclass[11pt,a4paper]{article}

% ============ PACKAGES ============
\usepackage[utf8]{inputenc}
\usepackage[T1]{fontenc}
\usepackage[french]{babel}

% Page geometry for A4 with better margins
\usepackage[
  a4paper,
  margin=2.5cm,
  top=3cm,
  bottom=3cm,
  includeheadfoot
]{geometry}

% Typography and fonts
\usepackage{lmodern}
\usepackage{microtype}
\usepackage{setspace}
\setstretch{1.15}

% Graphics and colors
\usepackage{graphicx}
\usepackage[dvipsnames,table]{xcolor}
\usepackage{tikz}

% Tables
\usepackage{booktabs}
\usepackage{tabularx}
\usepackage{multirow}
\usepackage{array}
\usepackage{longtable}

% Lists and formatting
\usepackage{enumitem}
\usepackage{parskip}

% Boxes
\usepackage{tcolorbox}

% Headers and footers
\usepackage{fancyhdr}

% Hyperlinks
\usepackage{hyperref}
\hypersetup{
  colorlinks=true,
  linkcolor=blue!50!black,
  urlcolor=blue!50!black,
  citecolor=blue!50!black
}

% ============ CUSTOM COLORS ============
\definecolor{maincolor}{RGB}{0,82,147}
\definecolor{secondcolor}{RGB}{0,120,190}
\definecolor{accentcolor}{RGB}{230,126,34}
\definecolor{lightgray}{RGB}{245,245,245}
\definecolor{tableheader}{RGB}{41,128,185}

% ============ CUSTOM COMMANDS ============
\newcommand{\sectionbox}[1]{%
  \vspace{0.5em}
  \noindent\colorbox{maincolor}{\parbox{\dimexpr\textwidth-2\fboxsep}{%
    \color{white}\Large\bfseries\centering #1%
  }}%
  \vspace{0.5em}
}

\newcommand{\subsectionbox}[1]{%
  \vspace{0.3em}
  \noindent\colorbox{secondcolor}{\parbox{\dimexpr\textwidth-2\fboxsep}{%
    \color{white}\large\bfseries #1%
  }}%
  \vspace{0.3em}
}

% ============ PAGE STYLE ============
\pagestyle{fancy}
\fancyhf{}
\fancyhead[L]{\small Rapport de suivi - Lean Manufacturing}
\fancyhead[R]{\small Mehdi Boumazzourh}
\fancyfoot[C]{\thepage}
\renewcommand{\headrulewidth}{0.4pt}
\renewcommand{\footrulewidth}{0.4pt}

% ============ DOCUMENT START ============
\begin{document}

% ============ TITLE PAGE ============
\begin{titlepage}
  \centering
  \vspace*{2cm}
  
  % Logo placeholders
  \begin{tabular}{p{0.3\textwidth}p{0.4\textwidth}p{0.3\textwidth}}
    \centering [Logo 1] & 
    \centering\textbf{\Large ECOLE SUPÉRIEURE DES INDUSTRIES\\DU TEXTILE ET DE L'HABILLEMENT} & 
    \centering [Logo 2]
  \end{tabular}
  
  \vspace{2cm}
  
  {\Huge\bfseries\color{maincolor} Rapport de suivi\par}
  
  \vspace{1.5cm}
  
  \begin{tcolorbox}[colback=lightgray,colframe=maincolor,width=0.8\textwidth,arc=3mm,boxrule=1.5pt]
    \centering
    \textbf{\Large École supérieure des industries textile et habillement}
    
    \vspace{0.5cm}
    
    \begin{tabular}{rl}
      \textbf{Année scolaire :} & 2025 - 2026 \\[0.3cm]
      \textbf{Nom et prénom de l'étudiant :} & Mehdi Boumazzourh \\[0.3cm]
      \textbf{Encadrant de l'ESITH :} & Mr SAMIR TETOUANI \\[0.3cm]
      \textbf{Domaine :} & Lean manufacturing et excellence opérationnelle \\[0.3cm]
      \textbf{Cycle :} & Génie industriel \\[0.3cm]
      \textbf{Filière :} & Chef de produit \\[0.3cm]
      \textbf{Niveau :} & 3\textsuperscript{ème} année \\
    \end{tabular}
  \end{tcolorbox}
  
  \vfill
  
  % Logo placeholder at bottom
  \includegraphics[width=0.3\textwidth]{example-image-a}
  
  \vspace{1cm}
  
\end{titlepage}

% ============ TABLE OF CONTENTS ============
\tableofcontents
\newpage

% ============ MAIN CONTENT ============

\section{Définition des termes généraux}

\subsectionbox{Introduction au Lean}

Le \textbf{Lean} (qui signifie « maigre » ou « sans superflu ») est une philosophie de gestion inspirée du système de production de \textbf{Toyota} (TPS). Son objectif principal est de maximiser la valeur pour le client tout en minimisant les gaspillages.

Bien que les deux termes soient souvent utilisés de manière interchangeable, ils désignent des périmètres légèrement différents :

\subsection{Le Lean Manufacturing (La dimension opérationnelle)}

C'est l'application historique du Lean au sein des \textbf{usines et des lignes de production}. Il se concentre sur l'optimisation des flux physiques.

\begin{itemize}[leftmargin=*,itemsep=0.3em]
  \item \textbf{Objectif :} Produire « au plus juste » (juste ce qu'il faut, quand il faut).
  \item \textbf{Action :} Éliminer les \textbf{Muda} (gaspillages) dans l'atelier, comme les stocks excessifs, les déplacements inutiles ou les défauts de fabrication.
  \item \textbf{Outils clés :} 5S (organisation), Kanban (flux tirés), SMED (réduction des temps de changement de série), TPM (maintenance).
\end{itemize}

\subsection{Le Lean Management (La dimension organisationnelle)}

Il s'agit de l'évolution du Lean Manufacturing à \textbf{l'ensemble de l'entreprise} (bureaux, RH, logistique, service client, direction). C'est autant un système de gestion qu'une culture d'entreprise.

\begin{itemize}[leftmargin=*,itemsep=0.3em]
  \item \textbf{Objectif :} Créer une culture d'amélioration continue (\textbf{Kaizen}) impliquant tous les collaborateurs.
  \item \textbf{Action :} Optimiser les processus administratifs et décisionnels, et surtout, développer les compétences des employés pour qu'ils résolvent eux-mêmes les problèmes.
  \item \textbf{Focus :} Le leadership, l'écoute du client et l'intelligence collective.
\end{itemize}

\subsection{Tableau Comparatif Rapide}

\begin{table}[H]
\centering
\begin{tabular}{@{}>{\bfseries}l>{\raggedright\arraybackslash}p{0.35\textwidth}>{\raggedright\arraybackslash}p{0.35\textwidth}@{}}
\toprule
\rowcolor{tableheader}
\color{white}\textbf{Caractéristique} & \color{white}\textbf{Lean Manufacturing} & \color{white}\textbf{Lean Management} \\
\midrule
Périmètre & L'atelier / La production & L'organisation entière \\
\rowcolor{lightgray}
Cible & Les machines, les flux, les pièces & Les personnes, les processus, la stratégie \\
Outils types & VSM, 5S, Kanban, Jidoka & Management Visuel, Kaizen, A3, Hoshin Kanri \\
\rowcolor{lightgray}
Résultats visés & Gain de productivité, réduction de stock & Agilité, satisfaction client, engagement salarié \\
\bottomrule
\end{tabular}
\caption{Comparaison Lean Manufacturing vs Lean Management}
\end{table}

\newpage
\section{Les 5 principes fondamentaux du Lean}

Pour réussir dans l'une ou l'autre de ces approches, on suit généralement ces 5 étapes :

\begin{enumerate}[leftmargin=*,itemsep=0.5em]
  \item \textbf{Définir la Valeur :} Qu'est-ce que le client est réellement prêt à payer ?
  
  \item \textbf{Identifier le Flux de Valeur (VSM) :} Cartographier toutes les étapes pour repérer celles qui ne servent à rien.
  
  \item \textbf{Créer un Flux Continu :} Faire en sorte que le produit/service avance sans interruption (ni attente, ni stock).
  
  \item \textbf{Mettre en place le Flux Tiré :} Ne produire que lorsqu'il y a une demande réelle du client.
  
  \item \textbf{Viser la Perfection :} Répéter le cycle pour s'améliorer sans cesse (Kaizen).
\end{enumerate}

\subsection{Comparaison entre excellence opérationnelle et lean management}

\begin{table}[H]
\centering
\small
\begin{tabular}{@{}>{\bfseries}l>{\raggedright\arraybackslash}p{0.35\textwidth}>{\raggedright\arraybackslash}p{0.35\textwidth}@{}}
\toprule
\rowcolor{tableheader}
\color{white}\textbf{Caractéristique} & \color{white}\textbf{Lean Management} & \color{white}\textbf{Excellence Opérationnelle} \\
\midrule
Définition & Philosophie visant à éliminer les gaspillages pour maximiser la valeur client. & Démarche systémique visant une performance durable et une culture d'amélioration. \\
\rowcolor{lightgray}
Focus principal & Flux, vitesse et réduction des gaspillages (Muda). & Culture, leadership, stratégie et résultats globaux. \\
Boîte à outils & 5S, VSM, Kanban, Kaizen, SMED. & Lean + Six Sigma + TPM + Management stratégique + Digital. \\
\rowcolor{lightgray}
Vision & « Comment faire mieux avec moins ? » & « Comment ancrer la performance dans l'ADN de l'entreprise ? » \\
Acteurs & Souvent perçu comme une démarche terrain (Gemba). & Implication totale, de la direction générale aux opérateurs. \\
\bottomrule
\end{tabular}
\caption{Lean Management vs Excellence Opérationnelle}
\end{table}

\newpage
\section{Définition des termes et de quelques méthodes}

\subsection{Définition des Termes et Méthodes Clés}

Avant d'aborder les problèmes, clarifions le vocabulaire technique :

\begin{description}[leftmargin=2cm,style=nextline,itemsep=0.5em]
  \item[\textbf{Lean Manufacturing}] Philosophie visant à éliminer les \textbf{Muda} (gaspillages) pour maximiser la valeur ajoutée avec le moins de ressources possible.
  
  \item[\textbf{Six Sigma}] Méthodologie statistique visant à réduire la \textbf{variabilité} des processus pour atteindre un taux de défaut quasi nul (3,4 défauts par million d'opportunités).
  
  \item[\textbf{Lean Six Sigma (LSS)}] L'hybride des deux : la vitesse et l'agilité du Lean combinées à la rigueur statistique du Six Sigma.
  
  \item[\textbf{Just-In-Time (JIT)}] Produire exactement ce qui est nécessaire, quand c'est nécessaire et en quantité voulue. Repose sur le flux tiré (\textbf{Kanban}).
  
  \item[\textbf{Kaizen}] Démarche d'amélioration continue par « petits pas », impliquant tout le personnel, du terrain à la direction.
  
  \item[\textbf{VSM (Value Stream Mapping)}] Cartographie du flux de valeur permettant de visualiser les gaspillages entre le fournisseur et le client.
\end{description}

\subsection{Typologie des Problèmes Industriels}

En industrie, on classe généralement les problèmes selon l'indicateur qu'ils dégradent (modèle \textbf{SQCDP}) :

\begin{enumerate}[leftmargin=*,itemsep=0.5em]
  \item \textbf{Qualité (Quality) :} Taux de rebuts trop élevé, non-conformités, réclamations clients, variabilité dimensionnelle.
  
  \item \textbf{Productivité/Coût (Cost) :} Temps de cycle trop long, stocks excessifs (WIP), pannes machines fréquentes, sous-optimisation de la main-d'œuvre.
  
  \item \textbf{Délai (Delivery) :} Retards de livraison, goulots d'étranglement dans le flux, manque de composants (ruptures).
  
  \item \textbf{Sécurité \& Environnement (Safety/Environment) :} Risques d'accidents, ergonomie des postes (TMS), consommation énergétique excessive.
\end{enumerate}

\subsection{Méthodologie DMAIC}

La méthodologie DMAIC est un processus structuré en 5 phases :

\begin{itemize}[leftmargin=*,itemsep=0.3em]
  \item \textbf{Define (Définir)} - Définir le problème et les objectifs
  \item \textbf{Measure (Mesurer)} - Mesurer la performance actuelle
  \item \textbf{Analyze (Analyser)} - Analyser les causes racines
  \item \textbf{Improve (Améliorer)} - Mettre en œuvre les solutions
  \item \textbf{Control (Contrôler)} - Contrôler et pérenniser les gains
\end{itemize}

\newpage
\subsection{8D vs DMAIC : Deux approches distinctes}

En milieu industriel, on distingue ces deux situations par la nature de l'\textbf{écart} (le « Gap ») :

\subsubsection{La 8D : La gestion de l'Anomalie (Restauration)}

\begin{itemize}[leftmargin=*,itemsep=0.3em]
  \item \textbf{Situation :} Le standard existe, il était maîtrisé, mais « quelque chose » a cassé le processus.
  \item \textbf{L'objectif :} C'est un retour vers le passé (revenir à l'état stable que l'on connaissait).
  \item \textbf{Le mot-clé :} Réactivité. On parle de « Correctif ».
  \item \textbf{Analogie :} Un athlète qui court habituellement le 100m en 10 secondes (objectif atteint) et qui, suite à une blessure, tombe à 15 secondes. La 8D, c'est le protocole de soin pour qu'il court à nouveau en 10 secondes.
\end{itemize}

\subsubsection{Le DMAIC : La gestion du Projet (Rupture)}

\begin{itemize}[leftmargin=*,itemsep=0.3em]
  \item \textbf{Situation :} Le processus est stable, il n'y a pas de « panne », mais le niveau actuel ne suffit plus pour être compétitif.
  \item \textbf{L'objectif :} C'est un bond vers le futur (atteindre un niveau de performance jamais atteint auparavant).
  \item \textbf{Le mot-clé :} Proactif. On parle d'Amélioration de rupture (Breakthrough).
  \item \textbf{Analogie :} Le même athlète court toujours en 10 secondes (stable), mais pour gagner les JO, il doit passer à 9,7 secondes. Il ne s'agit pas de soigner une blessure, mais de changer de régime, de technique de course et d'équipement. C'est le DMAIC.
\end{itemize}

\begin{table}[H]
\centering
\begin{tabular}{@{}lll@{}}
\toprule
\rowcolor{tableheader}
\color{white}\textbf{Situation} & \color{white}\textbf{Question à se poser} & \color{white}\textbf{Méthode} \\
\midrule
Chute soudaine & Qu'est-ce qui a changé par rapport à hier ? & 8D \\
\rowcolor{lightgray}
Objectif futur & Comment faire mieux que ce que l'on sait déjà faire ? & DMAIC \\
\bottomrule
\end{tabular}
\caption{Choix entre 8D et DMAIC}
\end{table}

\newpage
\section{Étude de cas : Projet Robotica}

\subsectionbox{Phase Définir (Define)}

\subsection{Problématique}

\begin{table}[H]
\small
\begin{tabularx}{\textwidth}{@{}lX@{}}
\toprule
\rowcolor{tableheader}
\multicolumn{2}{c}{\color{white}\textbf{Détails du Projet Robotica}} \\
\midrule
\textbf{Situation actuelle} & 
\begin{itemize}[nosep,leftmargin=1em]
  \item Mécontentement des donneurs d'ordres : retards de livraison et coût élevé des défauts qualité
  \item Menace de changement de fournisseur sous 6 mois si la situation n'est pas redressée
  \item Performance technique : 45 pannes/mois, temps de changement de référence de 3h, délai de livraison moyen de 7 jours
\end{itemize} \\
\rowcolor{lightgray}
\textbf{Situation future} & 
\begin{itemize}[nosep,leftmargin=1em]
  \item Clients satisfaits et fidélisés grâce à une meilleure fiabilité
  \item Délais de livraison réduits de 20\% (cible : 5,6 jours)
  \item Entreprise rentable avec une marge supérieure à 20\%
  \item Processus stabilisé avec une réduction significative des pannes et du temps SMED
\end{itemize} \\
\textbf{Objectifs} & 
\begin{itemize}[nosep,leftmargin=1em]
  \item Réduction du coût de revient de 15\%
  \item Réduction des délais de livraison de 20\%
  \item Augmentation de la marge à plus de 20\%
\end{itemize} \\
\rowcolor{lightgray}
\textbf{Équipe de projet} & 
\begin{itemize}[nosep,leftmargin=1em]
  \item Sponsor : Directeur Général (DG) de Robotica
  \item Leader : Responsable Excellence Opérationnelle (EO)
  \item Partenaires : Équipes R\&D (Allemagne) et Production/Montage (France)
\end{itemize} \\
\textbf{Macro planning} & 
\begin{itemize}[nosep,leftmargin=1em]
  \item Remise du rapport de diagnostic : sous 20 heures
  \item Phase de redressement globale : maximum 6 mois
\end{itemize} \\
\rowcolor{lightgray}
\textbf{Risques du projet} & 
\begin{itemize}[nosep,leftmargin=1em]
  \item Perte définitive des clients majeurs si le délai de 6 mois n'est pas respecté
  \item Instabilité critique de la production due au nombre élevé de pannes (45/mois)
  \item Inflexibilité face à la variation de la demande (60 à 100 commandes)
\end{itemize} \\
\textbf{Plan de communication} & 
\begin{itemize}[nosep,leftmargin=1em]
  \item Reporting direct et immédiat au DG (diagnostic sous 20h)
  \item Coordination inter-sites entre la R\&D allemande et l'usine française
  \item Suivi mensuel des indicateurs clés (KPI) de performance pour les donneurs d'ordres
\end{itemize} \\
\bottomrule
\end{tabularx}
\caption{Charte de projet DMAIC - Robotica}
\end{table}

\newpage
\subsectionbox{Phase Mesurer (Measure)}

\subsection{Mesure des KPI's}

\textbf{Voix du client (CTQ - Critical To Quality) :}
\begin{itemize}[leftmargin=*,itemsep=0.3em]
  \item KPI délai de livraison
  \item Taux de conformité qualité
  \item Coût total
\end{itemize}

\textbf{Voix du processus :}
\begin{itemize}[leftmargin=*,itemsep=0.3em]
  \item Taux de satisfaction client
  \item Manufacturing Lead Time (MLT)
  \item Ratio TVA/MLT
\end{itemize}

\subsection{Synthèse Globale}

Le total du Chiffre d'Affaires (CA) est de \textbf{60 386 360 €} pour une Quantité Totale vendue de \textbf{528 161 unités}, réparties sur \textbf{1 456 transactions}. La base de données couvre \textbf{10 produits uniques} et \textbf{5 clients uniques}.

\subsubsection{Prise de Vue Générale et Concentration des Revenus}

L'analyse révèle une forte concentration du Chiffre d'Affaires, tant au niveau des produits qu'au niveau des clients.

\paragraph{Concentration des Produits :}
\begin{itemize}[leftmargin=*,itemsep=0.3em]
  \item Les \textbf{5 premiers produits} (sur 10) génèrent \textbf{67.66\%} du Chiffre d'Affaires total
  \item Les produits \textbf{P02} et \textbf{P10} sont les principaux moteurs de revenus, suivis par \textbf{P04}, \textbf{P07}, et \textbf{P08}
  \item Les produits \textbf{P02} et \textbf{P10} sont également les produits les plus vendus en volume (quantité), avec des volumes très élevés par rapport aux autres produits du Top 5
\end{itemize}

\paragraph{Concentration des Clients :}
\begin{itemize}[leftmargin=*,itemsep=0.3em]
  \item Avec seulement 5 clients uniques, le top 5 des clients génère logiquement \textbf{100.00\%} du CA total
  \item Les trois principaux clients sont \textbf{Stäubli}, \textbf{Kawasaki} et \textbf{ABB}
\end{itemize}

\subsection{Implications Stratégiques}

D'un point de vue de gestionnaire industriel ou de data analyst, cette répartition impose des stratégies différenciées :

\subsubsection{Gestion des Stocks}
Pour \textbf{P02} et \textbf{P10}, une rupture de stock serait catastrophique. Ils nécessitent un suivi rigoureux (inventaire permanent, stock de sécurité optimisé). Pour les autres, une gestion sur commande ou un stock minimal suffit.

\subsubsection{Optimisation de la Production}
Il est probable que les lignes de production soient quasi exclusivement configurées pour P02 et P10. Le passage à la production des autres références génère probablement des « temps de changement de série » (SMED) coûteux par rapport au volume produit.

\subsubsection{Analyse de Rentabilité}
Il serait intéressant de vérifier si ces deux produits à fort volume sont aussi ceux qui génèrent le plus de marge, ou s'ils servent de « produits d'appel ».

\subsubsection{Stratégie Volume vs. Valeur}
\begin{itemize}[leftmargin=*,itemsep=0.3em]
  \item \textbf{P02 et P10} sont les produits de \textbf{flux} (grande consommation, faible prix unitaire, rotation rapide)
  \item \textbf{P04, P07 et P08} sont les produits de \textbf{spécialité} (haute valeur unitaire, probablement plus complexes à produire)
\end{itemize}

\subsubsection{Analyse du Prix Unitaire}
On constate un écart massif. Par exemple, le produit \textbf{P01} (2,7 millions de CA pour seulement 107 unités) a un prix unitaire environ \textbf{650 fois plus élevé} que le produit \textbf{P02}.

\subsubsection{Analyse Pareto}
Concentrez les efforts d'optimisation des coûts et de qualité sur les produits de la \textbf{Zone A} (P02, P10) et le début de la \textbf{Zone B} (P04, P07). Une réduction de coût de 5\% sur P02 a beaucoup plus d'impact sur la rentabilité globale qu'une réduction de 20\% sur P01.

\newpage
\subsection{Manufacturing Lead Time (MLT)}

Le Manufacturing Lead Time est décomposé comme suit :

\begin{equation}
\text{MLT} = \text{Temps machine} + \text{Temps de changement} + \text{Temps en-cours} + \text{Temps attente} + \text{Temps déplacement}
\end{equation}

Où :
\begin{itemize}[leftmargin=*,itemsep=0.3em]
  \item \textcolor{blue}{\textbf{Temps à valeur ajoutée (TVA)}} : Temps machine + Temps de changement de format
  \item \textcolor{red}{\textbf{Temps à non-valeur ajoutée (TNVA)}} : Temps des en-cours + Temps attente + Temps déplacement
\end{itemize}

\textbf{Ratio d'efficacité Lean :}
\begin{equation}
\text{Ratio} = \frac{\text{TVA}}{\text{MLT}} = \frac{\text{TVA}}{\text{TVA} + \text{TNVA}}
\end{equation}

\begin{itemize}[leftmargin=*,itemsep=0.3em]
  \item Si Ratio $\approx$ 1 : \textcolor{green}{\textbf{Application LEAN efficace}}
  \item Si Ratio $<$ 1 : \textcolor{red}{\textbf{TNVA trop important - Amélioration nécessaire}}
\end{itemize}

\newpage
\section{Conclusion et Recommandations}

Les analyses menées dans le cadre de ce projet DMAIC sur le cas Robotica ont permis d'identifier plusieurs axes d'amélioration prioritaires :

\begin{enumerate}[leftmargin=*,itemsep=0.5em]
  \item \textbf{Réduction des pannes machines} : De 45 pannes/mois à un niveau acceptable pour stabiliser la production
  
  \item \textbf{Optimisation du SMED} : Réduire les temps de changement de série de 3h à moins de 1h
  
  \item \textbf{Focus sur les produits stratégiques} : Concentrer les efforts d'amélioration sur P02 et P10 (Zone A de l'analyse Pareto)
  
  \item \textbf{Amélioration du MLT} : Réduire les temps à non-valeur ajoutée pour améliorer le ratio TVA/MLT
  
  \item \textbf{Satisfaction client} : Atteindre les objectifs de délai (5,6 jours) et de marge (>20\%)
\end{enumerate}

\vspace{1cm}

\begin{center}
\textit{Ce rapport constitue la base du diagnostic et des actions à mener\\
dans le cadre du projet d'excellence opérationnelle.}
\end{center}

\end{document}
